\subsection{Counterfactuals and Policy Evaluation}

Policy learning and evaluation combine prediction with causal questions about actions not actually taken.

\begin{definition} \label{def:sl-potential-outcomes} \hlt{Potential Outcomes}\\
For treatment $A\in\{0,1\}$, the potential outcomes $Y(1)$ and $Y(0)$ denote the outcomes that would occur under each treatment level.  The individual causal effect is $Y(1)-Y(0)$, while average effects take expectations over a population.
\end{definition}

\begin{definition} \label{def:sl-ignorability-propensity} \hlt{Ignorability and Propensity}\\
Treatment assignment is ignorable given covariates $X$ if $(Y(0),Y(1))$ is conditionally independent of $A$ given $X$.  The propensity score is $e(X)=P(A=1\mid X)$.
\end{definition}

\begin{definition} \label{def:sl-off-policy-evaluation} \hlt{Off-Policy Evaluation}\\
Off-policy evaluation estimates the value of a target policy using data collected under another policy.  Importance weighting corrects the action distribution when the logging probabilities are known and overlap holds.
\end{definition}

\begin{remark} \label{rem:sl-feedback-policy-learning} \hlt{Feedback}\\
When predictions influence actions, the data-generating process changes.  This connects causal learning, reinforcement learning, and deployment monitoring.
\end{remark}
